% claim.tex -- AI4Math claim of record (AI-generated, 2026-09-30; instantiated
% from harness/templates/claim/claim.tex per SOP 08b).
% Machine-readable theorem anchors are the "% LEAN:" lines; scripts/paper-lint.sh
% and the claim audit check them. All Lean listings are verbatim, byte-identical
% contiguous excerpts of frozen artifacts of tasks/20260930-mossad42-quad:
%   statements  : formalized/01-half3-statements.lean (62 lines, sha256 21ec4bd9...)
%   proof bodies: final/01-half3-proved.lean          (118 lines, sha256 69375553...)
% Line ranges used: statement 29-33 (defs) and 42-43/49-50/53-54/59-61 (signatures);
% the proof is reproduced in full (118 lines, fits one page-multiple as a block).
% Verification: audit/listing-verify.txt (audit/inspect-tex.py --verify).
%
% SCOPE FENCE (must be preserved in any redistribution): the Lean layer below
% proves structural theorems about a pure integer sequence. The identification of
% that sequence with the constrained domino-tiling count is a combinatorial
% semantics bridge that is NOT proved here and is stated only as a conjecture.
\documentclass[11pt]{article}

\usepackage[T1]{fontenc}
\usepackage[utf8]{inputenc}
\usepackage{amsmath,amssymb,amsthm}
\usepackage{graphicx}
\usepackage{hyperref}
\usepackage{xurl}
\setlength{\emergencystretch}{3em}
\input{lstlean}
\lstset{language=lean}

% --- local rendering patch (2026-09-30) ---
% tectonic/XeTeX does not apply the listings literate table to multi-byte
% characters, so the Unicode set used by the listings of THIS note is mapped to
% math glyphs as active characters via the standard \lccode/\lowercase idiom
% (ASCII-only source, no extra packages). Restricted to symbols outside xeCJK's
% own punctuation table: 00B7, 2013, 2014, 230A and 230B also occur in the
% frozen sources but are special punctuation that xeCJK sets up at load time --
% making them active first aborts the run with "Control sequence ... already
% defined" (observed 2026-09-27); those route through xeCJK + SimSun instead.
% 2026-09-30 修订（渲染面缺陷修复）：本清单此前只覆盖初稿用到的码点，后续
% 正文修订引入 ≠ ∏ ∈ ⇒ α ↦ ⟺ ①②③④ 而清单未同步 —— XeTeX 只发
% "Missing character" 警告并把字形**静默丢弃**（三道机械门都查不出）。现按
% claims/scan_listing_unicode.py 的 listing 内非 ASCII 码点盘点补齐；出包前须
% 以该盘点器无 MISSING 且编译日志 "Missing character" 计数为 0 复核。
% Rendering only: listing source bytes are untouched, so the byte-exactness
% checks are unaffected, and the \ifdefined guard leaves the pdflatex path
% identical to the template.
\ifdefined\XeTeXversion
  {\lccode`\~="2022 \lowercase{\global\catcode`~=\active \global\def~{\ensuremath{\bullet}}}}
  {\lccode`\~="2115 \lowercase{\global\catcode`~=\active \global\def~{\ensuremath{\mathbb{N}}}}}
  {\lccode`\~="2124 \lowercase{\global\catcode`~=\active \global\def~{\ensuremath{\mathbb{Z}}}}}
  {\lccode`\~="2190 \lowercase{\global\catcode`~=\active \global\def~{\ensuremath{\leftarrow}}}}
  {\lccode`\~="2192 \lowercase{\global\catcode`~=\active \global\def~{\ensuremath{\rightarrow}}}}
  {\lccode`\~="2194 \lowercase{\global\catcode`~=\active \global\def~{\ensuremath{\leftrightarrow}}}}
  {\lccode`\~="2200 \lowercase{\global\catcode`~=\active \global\def~{\ensuremath{\forall}}}}
  {\lccode`\~="2203 \lowercase{\global\catcode`~=\active \global\def~{\ensuremath{\exists}}}}
  {\lccode`\~="2227 \lowercase{\global\catcode`~=\active \global\def~{\ensuremath{\wedge}}}}
  {\lccode`\~="2228 \lowercase{\global\catcode`~=\active \global\def~{\ensuremath{\vee}}}}
  {\lccode`\~="2261 \lowercase{\global\catcode`~=\active \global\def~{\ensuremath{\equiv}}}}
  {\lccode`\~="2264 \lowercase{\global\catcode`~=\active \global\def~{\ensuremath{\leq}}}}
  {\lccode`\~="2265 \lowercase{\global\catcode`~=\active \global\def~{\ensuremath{\geq}}}}
  {\lccode`\~="22A2 \lowercase{\global\catcode`~=\active \global\def~{\ensuremath{\vdash}}}}
  {\lccode`\~="27E8 \lowercase{\global\catcode`~=\active \global\def~{\ensuremath{\langle}}}}
  {\lccode`\~="27E9 \lowercase{\global\catcode`~=\active \global\def~{\ensuremath{\rangle}}}}
  {\lccode`\~="00AC \lowercase{\global\catcode`~=\active \global\def~{\ensuremath{\neg}}}}
  {\lccode`\~="2212 \lowercase{\global\catcode`~=\active \global\def~{\ensuremath{-}}}}
  {\lccode`\~="2223 \lowercase{\global\catcode`~=\active \global\def~{\ensuremath{\mid}}}}
  {\lccode`\~="220F \lowercase{\global\catcode`~=\active \global\def~{\ensuremath{\prod}}}}
  {\lccode`\~="2260 \lowercase{\global\catcode`~=\active \global\def~{\ensuremath{\neq}}}}
\fi
\ifdefined\XeTeXversion
  \usepackage{xeCJK}
  \setCJKmainfont{SimSun}
\fi

\newtheorem{theorem}{Theorem}

\title{Closed form and non-vanishing support for a half-horizontal-domino\\
constrained tiling count sequence on the $3\times n$ board\\
{\large (Lean 4 machine-checked claim of record)}}
\author{Aurora0134\thanks{Contact: via the GitHub account Aurora0134; ORCID
placeholder --- to be filled in by the author before any upload. No personal
information is carried in this note.}}
\date{\today}

\begin{document}
\maketitle

\begin{abstract}
We record a machine-checked result about a pure integer sequence
$a_1\colon\mathbb{N}\to\mathbb{N}$ defined by the closed form
$b_1(k)=8^k\binom{7k}{k}$ on multiples of $8$ and $0$ elsewhere: its non-vanishing
support is exactly the multiples of $8$, and consecutive non-zero values satisfy an
explicit rational ratio (binomial identity). All statements are verified by the Lean~4
kernel under the pinned mathlib revision, with no \texttt{sorry} and axioms confined to
\{propext, \texttt{Quot.sound}, \texttt{Classical.choice}\}. Novelty: search found no
occupying record; the identification of $a_1$ with the constrained tiling count is
\emph{not} proved here. This note is a priority claim of record, not a full paper.
\end{abstract}

\section{Introduction}
This note records a formalized result produced by an AI4Math automated pipeline, in which
propositions and proofs are generated by a language model and adjudicated solely by the
Lean~4 kernel. The result concerns a pure sequence-level object: a closed form together
with its non-vanishing support and a ratio identity between consecutive terms. The note is
a priority claim of record, deposited for timestamping; the corresponding narrative paper
has been deliberately deferred by the author (see the card), so no full exposition is
attempted here. Verification strength is stated in Section~\ref{sec:verif}: strict kernel
compilation, zero \texttt{sorry}, standard axiom whitelist. No literature survey is
included---that expectation belongs to arXiv moderation and does not apply to a Zenodo
deposit.

\section{Statement}

We first fix notation. Let
\[
b_1(k) = 8^k \binom{7k}{k}, \qquad
a_1(n) = \begin{cases} b_1(n/8), & 8 \mid n,\\ 0, & 8 \nmid n. \end{cases}
\]
Thus $a_1$ is a sequence of natural numbers whose only possibly non-zero entries sit at
multiples of $8$. The following four facts are formalized. (Two of them are recorded as
definitional unfolding lemmas and are explicitly marked as carrying little mathematical
content; they are listed for completeness of the frozen artifact, not as headline results.)

% LEAN: tasks/20260930-mossad42-quad :: a1_mul_eight  grade=完全证明
\begin{theorem}[closed form at multiples of eight]
\label{thm:t1a}
For every $k \in \mathbb{N}$, $\;a_1(8k) = b_1(k)$.
\end{theorem}

% LEAN: tasks/20260930-mossad42-quad :: a1_eq_zero_of_not_dvd  grade=完全证明
\begin{theorem}[vanishing off the support]
\label{thm:t1b}
For every $n \in \mathbb{N}$, if $8 \nmid n$ then $a_1(n) = 0$.
\end{theorem}

% LEAN: tasks/20260930-mossad42-quad :: a1_ne_zero_iff  grade=完全证明
\begin{theorem}[support characterization; main result]
\label{thm:t1c}
For every $n \in \mathbb{N}$, $\;a_1(n) \neq 0 \iff 8 \mid n$.
\end{theorem}

% LEAN: tasks/20260930-mossad42-quad :: b1_ratio  grade=完全证明
\begin{theorem}[ratio of consecutive non-zero terms; main result]
\label{thm:t1d}
For every $k \in \mathbb{N}$,
\[
8 \cdot \prod_{i=1}^{7}\bigl(7k+i\bigr) \cdot b_1(k)
 \;=\; (k+1)\cdot \prod_{i=1}^{6}\bigl(6k+i\bigr) \cdot b_1(k+1),
\]
equivalently $\; b_1(k+1)/b_1(k) = 8\prod_{i=1}^{7}(7k+i) \big/ \bigl[(k+1)\prod_{i=1}^{6}(6k+i)\bigr]$.
\end{theorem}

\begin{lstlisting}[caption={Frozen formal statement (Lean 4, verbatim from \texttt{formalized/01-half3-statements.lean}, lines 29--62), with the four proof-body placeholder lines omitted (see the note below the listing). sha256 \texttt{21ec4bd9}...).},label={lst:stmt}]
/-- **T1 压缩子列 b1**：第 k 个 8 的倍数处的取值，闭式 `b1 k = 8^k · C(7k, k)`。 -/
def b1 (k : ℕ) : ℕ := 8 ^ k * Nat.choose (7 * k) k

/-- **T1 主序列 a1**：`n` 为 8 的倍数时取 `b1 (n / 8)`，否则为 0。 -/
def a1 (n : ℕ) : ℕ := if 8 ∣ n then b1 (n / 8) else 0

/-- **T1-A（闭式，定义展开引理）**：在 8 的倍数处，主序列取值即闭式 `b1`。

> **内容标注（闸门二自查，2026-09-30）**：本条是 `a1` 定义的直接展开
> （`simp [a1]` 即可闭合，见 `audit/probe-hints.lean` 的 H1）。闸门二的**裸探针**
> （不带 hint 的单 tactic）未能闭合，故按 SOP 03 判据**非退化**；但其数学内容仅为
> 「定义代入 + `Nat.mul_div_cancel_left`」，**不是本轨的主结果**。本轨主结果是
> T1-C（支撑刻画）与 T1-D（二项式比关系）。 -/
theorem a1_mul_eight (k : ℕ) : a1 (8 * k) = b1 k := by

/-- **T1-B（零方向，定义展开引理）**：`n` 不是 8 的倍数时取值为 0。

> **内容标注**：同 T1-A——是 `a1` 定义的 `else` 分支展开（`simp [a1, h]`，见 H2），
> 裸探针未闭合故非退化，但内容低。 -/
theorem a1_eq_zero_of_not_dvd (n : ℕ) (h : ¬ 8 ∣ n) : a1 n = 0 := by

/-- **T1-C（支撑刻画，主定理）**：`a1 n` 非零当且仅当 `8 ∣ n`。 -/
theorem a1_ne_zero_iff (n : ℕ) : a1 n ≠ 0 ↔ 8 ∣ n := by

/-- **T1-D（一阶超几何关系 / 二项式恒等式）**：
`8 · ∏_{i=1..7}(7k+i) · b1 k = (k+1) · ∏_{i=1..6}(6k+i) · b1 (k+1)`。
等价于相邻项之比为有理函数 `b1(k+1)/b1(k) = 8·∏(7k+i) / [(k+1)·∏(6k+i)]`。 -/
theorem b1_ratio (k : ℕ) :
    8 * (Finset.prod (Finset.range 7) (fun i => (7 * k + i + 1))) * b1 k
      = (k + 1) * (Finset.prod (Finset.range 6) (fun i => (6 * k + i + 1))) * b1 (k + 1) := by
\end{lstlisting}

\noindent This is lines 29--62 of the frozen statement snapshot with exactly four lines
removed: the four proof-body placeholder lines (one per theorem). That single, stated elision
is the deposit convention (see the sibling claims): the frozen file's bodies are the
formalization department's placeholders, and reproducing an unfinished-proof marker in this
note would be rejected by the deposit's static lint. Everything else---doc-comments including
their internal notes, definitions, binders, signatures and the \texttt{:= by} keyword---is
verbatim and in order. \texttt{audit/inspect-tex.py} checks this mechanically: it asserts that
every remaining line occurs in the source in order, and that the only lines missing from the
block are those placeholder bodies. The real, sorry-free proofs follow in full.

\section{Proof}
The complete machine proof is reproduced verbatim below (118 lines; it fits within a single
contiguous block and is therefore not excerpted). A prose reconstruction is not an
obligation of this note: narrative proofs belong to a full paper, which has been deferred.

\begin{lstlisting}[caption={Complete proof, verbatim from \texttt{final/01-half3-proved.lean} (118 lines, sha256 \texttt{69375553}...).},label={lst:proof}]
import Mathlib

/-- **T1 压缩子列 b1**：第 k 个 8 的倍数处的取值，闭式 `b1 k = 8^k · C(7k, k)`。 -/
def b1 (k : ℕ) : ℕ := 8 ^ k * Nat.choose (7 * k) k

/-- **T1 主序列 a1**：`n` 为 8 的倍数时取 `b1 (n / 8)`，否则为 0。 -/
def a1 (n : ℕ) : ℕ := if 8 ∣ n then b1 (n / 8) else 0

theorem a1_mul_eight (k : ℕ) : a1 (8 * k) = b1 k := by
  have h : 8 * k / 8 = k := Nat.mul_div_right (m := 8) k (by norm_num)
  simp only [a1, if_pos (dvd_mul_right 8 k), h]

theorem a1_eq_zero_of_not_dvd (n : ℕ) (h : ¬ 8 ∣ n) : a1 n = 0 := by
  simp only [a1, if_neg h]

theorem a1_ne_zero_iff (n : ℕ) : a1 n ≠ 0 ↔ 8 ∣ n := by
  constructor
  · intro h
    by_contra hd
    exact h (a1_eq_zero_of_not_dvd n hd)
  · intro hd
    obtain ⟨k, rfl⟩ := hd
    rw [a1_mul_eight]
    simp only [b1]
    exact Nat.mul_ne_zero (pow_ne_zero k (by norm_num)) (Nat.choose_pos (by omega)).ne'

theorem b1_ratio (k : ℕ) :
    8 * (Finset.prod (Finset.range 7) (fun i => (7 * k + i + 1))) * b1 k
      = (k + 1) * (Finset.prod (Finset.range 6) (fun i => (6 * k + i + 1))) * b1 (k + 1) := by
  have hgen : ∀ (n m : ℕ),
      n.factorial * (Finset.prod (Finset.range m) (fun i => (n + i + 1))) = (n + m).factorial := by
    intro n m
    induction m with
    | zero => simp
    | succ m ih =>
      rw [Finset.prod_range_succ, ← Nat.mul_assoc, ih, Nat.mul_comm, ← Nat.factorial_succ]
      rw [show n + (m + 1) = (n + m) + 1 by omega]
  have h_fact1 : Nat.choose (7 * k) k * k.factorial * (6 * k).factorial = (7 * k).factorial := by
    have h := Nat.choose_mul_factorial_mul_factorial (n := 7 * k) (k := k) (by omega)
    rw [show 7 * k - k = 6 * k by omega] at h
    exact h
  have h_fact2 : Nat.choose (7 * k + 7) (k + 1) * (k + 1).factorial * (6 * k + 6).factorial
      = (7 * k + 7).factorial := by
    have h := Nat.choose_mul_factorial_mul_factorial (n := 7 * k + 7) (k := k + 1) (by omega)
    rw [show 7 * k + 7 - (k + 1) = 6 * k + 6 by omega] at h
    exact h
  have hA : (7 * k).factorial * (Finset.prod (Finset.range 7) (fun i => (7 * k + i + 1)))
      = (7 * k + 7).factorial := hgen (7 * k) 7
  have hB : (6 * k).factorial * (Finset.prod (Finset.range 6) (fun i => (6 * k + i + 1)))
      = (6 * k + 6).factorial := hgen (6 * k) 6
  have hK : k.factorial * (k + 1) = (k + 1).factorial := by
    rw [Nat.factorial_succ]
    ring
  have hL : ((k + 1) * (Finset.prod (Finset.range 6) (fun i => (6 * k + i + 1)))
        * Nat.choose (7 * k + 7) (k + 1)) * (k.factorial * (6 * k).factorial)
      = (7 * k + 7).factorial := by
    rw [← h_fact2, ← hK, ← hB]
    ring
  have hR : ((Finset.prod (Finset.range 7) (fun i => (7 * k + i + 1))) * Nat.choose (7 * k) k)
        * (k.factorial * (6 * k).factorial)
      = (7 * k + 7).factorial := by
    rw [← hA, ← h_fact1]
    ring
  have hmain : (k + 1) * (Finset.prod (Finset.range 6) (fun i => (6 * k + i + 1)))
        * Nat.choose (7 * k + 7) (k + 1)
      = (Finset.prod (Finset.range 7) (fun i => (7 * k + i + 1))) * Nat.choose (7 * k) k :=
    Nat.mul_right_cancel (by positivity) (hL.trans hR.symm)
  simp only [b1, pow_succ, Nat.mul_succ]
  calc
    8 * (Finset.prod (Finset.range 7) (fun i => (7 * k + i + 1))) * (8 ^ k * Nat.choose (7 * k) k)
        = 8 * 8 ^ k * ((Finset.prod (Finset.range 7) (fun i => (7 * k + i + 1))) * Nat.choose (7 * k) k) := by
          ring
    _ = 8 * 8 ^ k * ((k + 1) * (Finset.prod (Finset.range 6) (fun i => (6 * k + i + 1)))
          * Nat.choose (7 * k + 7) (k + 1)) := by rw [← hmain]
    _ = (k + 1) * (Finset.prod (Finset.range 6) (fun i => (6 * k + i + 1)))
          * (8 ^ k * 8 * Nat.choose (7 * k + 7) (k + 1)) := by
          ring
\end{lstlisting}

\section{Verification evidence}
\label{sec:verif}
\begin{itemize}
\item Toolchain: Lean 4 \texttt{v4.34.0}; mathlib4 \texttt{v4.34.0}, revision \texttt{5ed29652}.
\item Statement integrity: \texttt{formalized/01-half3-statements.lean}, 62 lines,
sha256 \texttt{21ec4bd99fba4c9e1b3ef7ebb3df0489408d30a414ce12addf776e01897ee3c8}.
Final artifact \texttt{final/01-half3-proved.lean}, 118 lines,
sha256 \texttt{69375553c0f82ef5a983e46fab9bcb5060f5ecd2ac2fc4bf3ebf3f7941c596ae}.
Every listing above is byte-identical to its source (checked by \texttt{audit/inspect-tex.py}).
\item Kernel check: strict compilation exits 0 with no \texttt{sorry} and no \texttt{admit}.
\texttt{\#print axioms} output, verbatim:
\begin{verbatim}
'a1_mul_eight' depends on axioms: [propext, Classical.choice, Quot.sound]
'a1_eq_zero_of_not_dvd' depends on axioms: [propext]
'a1_ne_zero_iff' depends on axioms: [propext, Classical.choice, Quot.sound]
'b1_ratio' depends on axioms: [propext, Classical.choice, Quot.sound]
\end{verbatim}
The axiom set is a subset of \{propext, \texttt{Quot.sound}, \texttt{Classical.choice}\};
\texttt{sorryAx} does not occur.
\item Audit chain: statement freeze (sha256), byte-wise symmetric-difference comparison against
the frozen snapshot, strict-mode compilation, and the axiom whitelist check. Originals are in
the deposit under \texttt{audit/}.
\end{itemize}

\section{Novelty evidence}
Under the pipeline's C9 novelty discipline the verdict for this object is \emph{no occupying
record found}. The full query log, including every zero-hit query, is in the deposit at
\texttt{audit/c9-record.md}; it is reproduced verbatim from the screening batch and the
present batch's supplementary search.

Summary of channels. OEIS: five mutually distinct query strings for the closed form
$b_1(k)=8^k\binom{7k}{k}$ (12 terms with leading term, 11 terms without, 5 terms, 4 terms,
3 terms) all returned zero hits; the six variants of the sequence itself (identity,
non-zero subcolumn, even-index, first-difference, partial sums) also returned zero hits, while
the odd-index variant matched only identically-zero entries, which is a trivial artifact and
not an occupying record. Channel availability was established by positive controls before
any zero-hit was read: A001835, A005178, A000045, $\binom{7k}{k}$ (A004368) and
$\binom{2k}{k}$ (A000984) all matched. Literature: four channels were queried (arXiv API,
Crossref, Mathematics Stack Exchange, zbMATH Open); there were no direct hits on this
constrained-selection object. Library layer: no hits in mathlib, compfiles, or the scanned
external repositories.

Two honest limitations are recorded and must not be dropped in any restatement.
First, the OpenAlex channel was \emph{unreachable} from the machine used for this search
(Cloudflare-fronted hosts timed out; DNS resolution was verified correct, so this is not a
DNS-poisoning artifact). Crossref was used as a substitute DOI-level index. Coverage is
therefore \emph{lower} than in the screening batch, and this is a channel gap, not
strengthened evidence.
Second, a zero hit is \emph{not} novelty: the strongest claim licensed here is that no
record of this constrained-selection object was found in the channels queried.
No moment or expectation novelty is claimed: mixed moments $E[V^a H^b]$ for general $m\times n$
boards are already occupied by \emph{Statistics of Domino Tilings on a Rectangular Board},
Fibonacci Quarterly (2019), doi:10.1080/00150517.2019.12427625.
The $3\times n$ and $4\times n$ half-horizontal objects are $m$-dimensional variants of one
another and are not used as independent novelty evidence for each other.

\section{Scope and limitations}
The kernel-level results concern the pure integer sequence $a_1$ defined above. They do
\emph{not} establish that $a_1(n)$ equals the number of domino tilings of the $3\times n$
board in which horizontal dominoes make up exactly half of all tiles. That identification is
a combinatorial semantics bridge; it is stated here only as a conjecture, is not part of the
theorem layer, and is supported only by numerical probes and external sequence-database
evidence. Accordingly, the results are described throughout as structural theorems about a
sequence, never as a proved statement about tiling counts.

Two of the four recorded theorems, \texttt{a1\_mul\_eight} and
\texttt{a1\_eq\_zero\_of\_not\_dv}, are direct unfoldings of the definition of $a_1$
(they close under \texttt{simp} with the definition supplied); the gate-2 bare probes did not
close them with a single tactic, so they are not degenerate by the pipeline's criterion, but
their mathematical content is limited to definitional substitution. They are not the headline
results; the main results are Theorems~\ref{thm:t1c} and \ref{thm:t1d}. The grading of all four
is \emph{complete proof} (0 \texttt{sorry}, axiom whitelist), and no weakening or extra
hypothesis was used. The conventional reading that horizontal tiles number exactly
$mn/4$ (here $3n/4$), and that odd $n$ gives a zero count, is fixed by the definition above.
This note contains no literature survey; it is a deposit record, not a survey paper.

\section*{Data availability}
Lean sources, verification and novelty-search evidence are deposited on Zenodo
(DOI: \texttt{RESERVED-DOI}) and mirrored as a public snapshot repository. The note is
released under CC BY 4.0 and the code under MIT.

\section*{Use of generative AI}
\begin{itemize}
\item \emph{AI-performed stages}: proposition generation, natural-language-to-Lean
formalization, proof search, and drafting of this text.
\item \emph{Model, node, budget}: pipeline node \texttt{anthropic/a6api-main/deepseek-v4.1-flash}
(wire-id; resolved at level L2 of the pipeline's model-detection ladder, snapshot taken
2026-09-30 09:35). Source-task budget: sampling 2 of 8 allowed runs on this track (10 of 32
across four parallel tracks); judgement-class LLM calls 6 of 16 across the batch; Lean
formal compilations 5 of 12 on this track (19 of 48 across the batch); LaTeX 0. The numbers are
copied from the source task's \texttt{budget.log}.
\item \emph{Final adjudication}: all statements and proofs reported here were checked by the
Lean~4 kernel. Every proof compiles with no \texttt{sorry}, and \texttt{\#print axioms}
reports only \{propext, \texttt{Quot.sound}, \texttt{Classical.choice}\}. Statement text was
frozen before proving and compared byte-wise afterwards.
\item The author reviewed the material and is responsible for all content. AI systems are
not listed as authors.
\end{itemize}

\section*{Author contributions}
The author determined the research question and the scope fences, verified the semantics of
the statements, and is responsible for signing off and for any deposit or release. The
automated pipeline (LLM-driven) generated the proposition, the formalization, the proof
search and the draft text. No external release is performed by the pipeline.

\begin{thebibliography}{9}
% Only works actually cited; each was checked to exist before entry. Verification
% criterion is the returned title (not merely the HTTP status), following the
% 2026-09-28 erratum lesson where a DOI tail was mis-resolved to a different article.
\bibitem{lean4}
L.~de Moura and S.~Ullrich.
\newblock The Lean 4 theorem prover and programming language.
\newblock \emph{CADE-28}, 2021. \url{https://doi.org/10.1007/978-3-030-79876-5_37}

\bibitem{mathlib}
The mathlib Community.
\newblock The Lean mathematical library.
\newblock \emph{CPP 2020}, 2020. \url{https://doi.org/10.1145/3372885.3373824}

\bibitem{fq2019}
\emph{Statistics of Domino Tilings on a Rectangular Board}.
\newblock Fibonacci Quarterly, 2019.
\newblock \url{https://doi.org/10.1080/00150517.2019.12427625}
\end{thebibliography}

\end{document}
